% !TeX root = ../../Book2.tex
\chapter{高维波动方程}
\label{b2:ch:11}
\BookChapterNumber{b2:ch:11}
\begin{chapterabstract}
高维波动方程同时具有球状传播、能量守恒与振荡抵消。本章由球面平均导出 Euler--Poisson--Darboux 方程，在三维恢复 Kirchhoff 公式，再用降维法解释二维尾部以及一般奇偶维的差异。分布基本解使球面核与球内核得到统一表述，Fourier 乘子则把演化延拓到能量空间。选读部分证明三维频率局部化后的逐点衰减，说明正则性代价以及进入时空估计的动机。
\end{chapterabstract}
\begin{learninggoals}
\begin{itemize}
\item\label{b2:item:11-goal-mean}由散度公式计算球面平均的半径导数，识别球心、半径与维数的不同作用。
\item\label{b2:item:11-goal-representation}从一维方程恢复三维 Kirchhoff 公式，并追踪降维时的面积权重。
\item\label{b2:item:11-goal-support}区分有限传播和强 Huygens 原理，说明一维及偶维尾部为何不矛盾。
\item\label{b2:item:11-goal-energy}把球面表示、延迟基本解和 Fourier 能量解联系起来，处理零频率及低正则初值。
\item\label{b2:item:11-goal-dispersion}在选读部分证明三维频带核估计，检验去掉频带条件后会失败的范数结论。
\end{itemize}
\end{learninggoals}

本章以\cref{b2:thm:09-dalembert}及\secref{b2:sec:0903}的能量极限为直接前置；球面积分和散度公式取自\chapref{b2:ch:01}，Fourier 方法使用\chapref{b2:ch:04}的酉性与\chapref{b2:ch:05}的分布语言。有限区间谱理论是另一条支线，本章不把\chapref{b2:ch:10}作为证明前置。

\ProseParagraphBreak

\secref{b2:sec:1101}建立球平均工具，\secref{b2:sec:1102}完成三维表示与局部能量，\secref{b2:sec:1103}讨论维数、基本解与能量空间。\secref{b2:sec:1104}可在首次阅读时略去。Fritz John 的《Partial Differential Equations》\cite{John1991Partial}适合伴读特征几何与波传播；Lawrence C. Evans 的《Partial Differential Equations》\cite{Evans2010Partial}可供进一步比较球平均、能量与低正则解的安排。两部书的具体归一化应与本章的波速和 \(2\pi\) 频率约定分别核对。

% !TeX root = ../../Book2.tex
\section{球平均法}
\label{b2:sec:1101}

一维公式依赖左右两个方向，高维空间却有连续无穷多个传播方向。对以观察点为球心的球面取平均，可以把这些方向组织成一个半径变量。得到的方程仍记得空间维数：球面的面积增长会在半径导数中产生一个系数。

\subsection{球面平均}
\label{b2:subsec:110101}

记 \(\sigma_{n-1}=\mathcal H^{n-1}(\mathbb S^{n-1})\) 为单位球面的面积；特别，\(\sigma_1=2\pi\)、\(\sigma_2=4\pi\)。对连续函数 \(h:\mathbb R^n\to\mathbb C\)，定义其\term[qiumian-pingjun]{球面平均}[spherical mean]为
\begin{equation}\label{b2:eq:11-spherical-mean}
 M_rh(x)\coloneqq\dfrac1{\sigma_{n-1}}
 \int_{\mathbb S^{n-1}}h(x+r\theta)\,\mathrm{d}\sigma(\theta),\quad r\in\mathbb R.
\end{equation}
积分测度是标准球面面积测度。换元 \(\theta\mapsto-\theta\) 表明平均关于 \(r\) 为偶函数；令 \(r=0\) 则得到 \(M_0h(x)=h(x)\)。若 \(r>0\)，亦可把它写成 \(\partial B_r(x)\) 上的面积积分除以 \(\sigma_{n-1}r^{n-1}\)。
\symidx{\(M_rh\)：半径为 \(r\) 的球面平均}

\ProseParagraphBreak

这里的球心 \(x\) 仍是自由变量。即使 \(h\) 不径向，每个球心附近也有自己的平均函数；不应把 \(M_rh(x)\) 误当作只依赖 \(r\) 的一个全局径向化结果。对平滑 \(h\)，求导可穿过紧球面上的积分，因此 \(\Delta_xM_rh=M_r\Delta h\)。另一方面，半径导数满足另一个表达式，正是它把空间算子转成单变量算子。

\begin{lemma}[球面平均的微分恒等式]\label{b2:lem:11-spherical-identity}
若 \(h\in C^2(\mathbb R^n)\)、\(n\geq2\)，则对 \(r>0\)，
\begin{equation}\label{b2:eq:11-epd-spatial}
 \partial_r^2M_rh+\dfrac{n-1}{r}\partial_rM_rh=M_r\Delta h=\Delta_xM_rh.
\end{equation}
当 \(r\to0\) 时，\(M_rh=h+r^2\Delta h/(2n)+o(r^2)\)，且 \(\partial_rM_rh|_{r=0}=0\)。
\end{lemma}
\begin{proof}
对\eqref{b2:eq:11-spherical-mean}求导，再换到半径为 \(r\) 的球面。由散度公式，
\[
 r^{n-1}\partial_rM_rh(x)
 =\dfrac1{\sigma_{n-1}}\int_{\partial B_r(x)}\partial_\nu h\,\mathrm{d}\sigma
 =\dfrac1{\sigma_{n-1}}\int_{B_r(x)}\Delta h(y)\,\mathrm{d}y.
\]
将球内积分写成半径积分，对 \(r\) 求导，右端为 \(r^{n-1}M_r\Delta h\)，除以 \(r^{n-1}\) 即得恒等式。

零点展开来自二阶 Taylor 公式。球面对称性给出 \(\int\theta_j\,\mathrm{d}\sigma=0\)，不同坐标的交叉项积分为零，而各 \(\theta_j^2\) 的积分相同，其总和为 \(\sigma_{n-1}\)。所以每项为 \(\sigma_{n-1}/n\)，得到二阶系数。Taylor 余项在紧球面方向上一致为 \(o(r^2)\)，可直接积分。第一半径导数为零也可由 \(h\) 的一阶展开得到。
\end{proof}

例如对 \(h(x)=|x|^2\)，有 \(M_rh(x)=|x|^2+r^2\)，左侧等于 \(2n\)，与 \(\Delta h\) 一致。若把球面平均误写成没有除面积的积分，额外的 \(r^{n-1}\) 将改变整个微分方程。

\subsection{Euler--Poisson--Darboux 方程}
\label{b2:subsec:110102}

设 \(u_{tt}=c^2\Delta u\) 为平滑解，固定观察中心 \(x\)，并记 \(U(t,r)=M_ru(t,\cdot)(x)\)。由前一引理，
\begin{equation}\label{b2:eq:11-epd}
 U_{tt}=c^2\left(U_{rr}+\dfrac{n-1}{r}U_r\right),\quad r>0.
\end{equation}
这是本章所用的\term[euler-poisson-darboux-fangcheng]{Euler--Poisson--Darboux 方程}[Euler--Poisson--Darboux equation]。系数在 \(r=0\) 看似奇异，但平均来自一个偶函数，且 \(U_r(t,0)=0\)。它在零点的行为不是可以任意添加的边界数据，而由球面平均的定义强制决定。

\ProseParagraphBreak

三维时系数为 \(2/r\)。乘以 \(r\)，令 \(V(t,r)=rU(t,r)\)，就有
\[
 V_{tt}=c^2V_{rr}.
\]
平均 \(U\) 为偶函数，故 \(V\) 为奇函数，并在 \(r=0\) 光滑穿过。若初始数据为 \((f,g)\)，则
\[
 V(0,r)=rM_rf(x),\quad V_t(0,r)=rM_rg(x).
\]
现在可在整条 \(r\) 轴上使用\cref{b2:thm:09-dalembert}。最后通过 \(u(t,x)=V_r(t,0)\) 恢复原解。这里的径向变量只服务于固定球心的平均，并未假设原始初值关于原点径向。

\subsection{降维}
\label{b2:subsec:110103}

把径向方程转成一维方程，与把一个高维解限制到较少变量，是两个不同步骤。后一种构造称为\term[jiangweifa]{降维法}[method of descent]：若高维初值不依赖最后一个坐标，且得到的解也具有这种独立性，那么最后一个坐标的二阶导数为零，方程就成为低一维的波方程。

\ProseParagraphBreak

该方法需要保留解的适用范围。把一个紧支集二维函数视作三维函数后，它沿第三坐标无限延伸，已经不属于三维 \(L^2\)。所以不能用只对三维有限总能量解成立的表示来完成降维；下一节的球面表示只需局部光滑初值，没有无穷远积分条件，因而适合这个目的。

\ProseParagraphBreak

几何上，三维球面的上下半球投影到二维圆盘，各自带来面积伸缩因子 \((1-|z|^2)^{-1/2}\)。它在圆周附近发散，却可积，因为
\[
 \int_{B_1(0)\subset\mathbb R^2}\dfrac{\mathrm{d}z}{\sqrt{1-|z|^2}}=2\pi.
\]
因此球面上的传播下降到平面以后，会变成整个圆盘上的加权传播。奇偶维差异的第一个例子已经出现在这一投影里，不能只把三维公式的球面维数改掉。


% !TeX root = ../../Book2.tex
\section{三维波方程}
\label{b2:sec:1102}

三维球平均乘以半径后满足一维波方程，这使观察点的位移由一个球面上的初值表示。球面半径随时间增长，但球内较早经过的点不会持续出现在齐次表示中。本节把这一精确的支集信息与一般能量方法所得的球内依赖区别开来。

\subsection{Kirchhoff 公式}
\label{b2:subsec:110201}

\begin{theorem}[Kirchhoff]\label{b2:thm:11-kirchhoff}
设 \(f\in C^3(\mathbb R^3)\)、\(g\in C^2(\mathbb R^3)\)。三维齐次初值问题 \(u_{tt}-c^2\Delta u=0\)、\(u(0)=f\)、\(u_t(0)=g\) 的经典解为
\begin{equation}\label{b2:eq:11-kirchhoff}
 u(t,x)=\partial_t\bigl(tM_{ct}f(x)\bigr)+tM_{ct}g(x).
\end{equation}
对 \(t>0\)，等价地，
\[
 u(t,x)=\partial_t\left(\dfrac1{4\pi c^2t}
       \int_{|y-x|=ct}f(y)\,\mathrm{d}\sigma(y)\right)
       +\dfrac1{4\pi c^2t}\int_{|y-x|=ct}g(y)\,\mathrm{d}\sigma(y).
\]
\end{theorem}
\begin{proof}
先对一个足够平滑的解求表示。按上一节的记号，令 \(H(r)=rM_rf(x)\)、\(J(r)=rM_rg(x)\)，两者均为奇函数。对 \(V=rU\) 使用一维公式并在 \(r=0\) 求导，得到
\[
 V_r(t,0)=H'(ct)+\dfrac1cJ(ct)
 =M_{ct}f(x)+ct\,\partial_rM_rf(x)|_{r=ct}+tM_{ct}g(x).
\]
右端正是\eqref{b2:eq:11-kirchhoff}。这个计算揭示了公式来源，存在性还需直接验证。

对任意 \(h\in C^2\)，令 \(W_h(t,x)=tM_{ct}h(x)\)。球面平均恒等式给出 \((W_h)_{tt}=c^2\Delta W_h\)，且 \(W_h(0)=0\)、\((W_h)_t(0)=h\)。于是 \(u=\partial_tW_f+W_g\) 满足波方程；\(f\in C^3\) 确保位移部分具有所需二阶导数。偶性及零点 Taylor 展开使 \(\partial_tW_f(0)=f\)、\(\partial_t^2W_f(0)=0\)，给出两项初值。球面平均化为面积积分产生 \(4\pi c^2t\) 的分母。经典解的唯一性由下一命题的局部能量论证得到，不附加无穷远条件。
\end{proof}

这称为\term[kirchhoff-gongshi]{Kirchhoff 公式}[Kirchhoff's formula]。在单位球面上展开时间导数，可把它写成
\[
 u(t,x)=\dfrac1{4\pi}\int_{\mathbb S^2}
 \bigl(f(x+ct\theta)+ct\,\theta\cdot\nabla f(x+ct\theta)
              +t\cdot g(x+ct\theta)\bigr)\,\mathrm{d}\sigma(\theta).
\]
因此“只依赖球面上的位移”须理解为还读取初始位移的法向导数；只知道两个函数在球面上的值相等，未必使它们的解相等。若数据在该球面的一个邻域相等，则所有所需数据当然都相等。

\ProseParagraphBreak

例如取 \(f(x)=|x|^2\)、\(g=0\)。由 \(M_{ct}f=|x|^2+c^2t^2\)，得到 \(u=|x|^2+3c^2t^2\)。它满足 \(u_{tt}=6c^2=c^2\Delta u\)。这份数据不具有有限总能量，却仍可使用局部球面公式。对空间常初速度 \(g=1\)、\(f=0\)，公式给出 \(u=t\)，再次核对了单位质量归一化。

\subsection{传播速度}
\label{b2:subsec:110202}

下面的计算适用于任何空间维数 \(n\)，也为随后各种表示公式提供统一的唯一性依据。令
\[
 e=\dfrac12(|u_t|^2+c^2|\nabla u|^2),\quad
 E_B(t)=\int_{B_{c(t_0-t)}(x_0)}e(t,x)\,\mathrm{d}x.
\]
局部守恒式为 \(\partial_te=c^2\operatorname{div}\operatorname{Re}(\overline{u_t}\nabla u)\)。球面随时间以法向速度 \(-c\) 收缩，因此出现一个符号有利的面积项。

\begin{proposition}[波方程的局部能量]\label{b2:prop:11-local-energy}
设 \(n\geq1\)，\(u\) 为收缩球锥邻域中的经典齐次解。则
\[
 E_B'(t)=\int_{\partial B_{c(t_0-t)}(x_0)}
 \bigl(c^2\operatorname{Re}(\partial_\nu u\,\overline{u_t})-c\cdot e\bigr)
 \,\mathrm{d}\sigma\leq0.
\]
若初始位移和速度在底球上为零，则 \(u\) 在整个锥内为零。因此任意维全空间经典初值问题至多有一个解，且其传播速度不超过 \(c\)。
\end{proposition}
\begin{proof}
对移动球积分求导并使用散度公式得到所示等式。将梯度分成法向分量和切向分量，积分被积函数等于
\[
 -\dfrac c2|u_t-c\partial_\nu u|^2
 -\dfrac{c^3}{2}|\nabla_{\mathrm{tan}}u|^2.
\]
所以局部能量不增。零初值使初始梯度和速度为零，因此能量恒为零。内部的 \(u_t\) 和 \(\nabla u\) 随之为零，沿留在锥内的竖直线积分到零时刻得到 \(u=0\)，边界由连续性补上。对两个解作差并让 \(x_0,t_0\) 任意，得到唯一性。比较局部相同的数据，同样给出球内依赖和有限传播。
\end{proof}

具有外力时，只在局部能量导数中增加 \(\operatorname{Re}\int F\overline{u_t}\)。若两个解的外力在过去锥内相等，其差仍满足上述齐次论证。因此一般外力的影响也沿速度不超过 \(c\) 的路径传递，不能从一个源点瞬时到达远处。

\subsection{Huygens 原理}
\label{b2:subsec:110203}

一般能量方法允许初始数据在整个底球上起作用；三维 Kirchhoff 公式却只读取球面上的数据及位移的一阶导数。这一更强的无尾传播性质称为\term[huygens-yuanli]{Huygens 原理}[Huygens' principle]，本章使用它的强形式：固定正时间的初速度传播核支撑在传播球面上，而非整个球内。

\begin{corollary}[三维无尾传播]\label{b2:cor:11-huygens}
若三维齐次初值 \(f,g\) 光滑并支撑在 \(\overline{B_R(0)}\) 内，则对 \(t>0\)，
\begin{equation}\label{b2:eq:11-shell}
 \operatorname{supp}u(t,\cdot)\subseteq
 \{\,x:\max\{0,ct-R\}\leq |x|\leq ct+R\,\}.
\end{equation}
特别，在固定位置 \(x\)，当 \(ct>|x|+R\) 时解为零。
\end{corollary}
\begin{proof}
若 \(|x|>ct+R\)，球面 \(\partial B_{ct}(x)\) 不与数据支集相交。若 \(ct>|x|+R\)，数据支集严格处于这个球的内部，也不与球面相交。在这两种情形中，球面都有一个邻域使 \(f,g\) 及 \(\nabla f\) 为零。展开的 Kirchhoff 公式遂为零。非零区域的闭包包含于所示闭壳层，便得到支集结论。
\end{proof}

内半径在 \(ct<R\) 时取零，不能写成负半径的球面。等号边界也只用于描述支集外界，不保证解在整个壳层上非零。如果外力持续存在，新产生的球面响应可以在内部叠加；无尾陈述首先针对齐次问题或单次瞬时扰动，不意味着持续振动源在身后必定留下空区。


% !TeX root = ../../Book2.tex
\section{高维情形}
\label{b2:sec:1103}

维数进入球平均的径向系数，也进入初速度传播核的支集。以下先把三维结果降到二维，再给出任意奇偶维的表示。最后改用 Fourier 变换构造有限能量解，使点值表示、分布基本解与范数演化得到同一个接口。

\subsection{奇偶维差异}
\label{b2:subsec:110301}

设 \(f,g\) 为 \(\mathbb R^2\) 上的光滑函数，把它们延拓为与第三坐标无关的函数。三维 Kirchhoff 公式显然也与第三坐标无关，因此给出二维解。单位球面的上下半球是 \((z,\pm\sqrt{1-|z|^2})\)，其面积元分别为 \((1-|z|^2)^{-1/2}\,\mathrm{d}z\)，故
\[
 M_{ct}^{(3)}\widetilde h(x,0)=\dfrac1{2\pi}
 \int_{|z|<1}\dfrac{h(x+ctz)}{\sqrt{1-|z|^2}}\,\mathrm{d}z.
\]
这里上标只用来区别球平均所处的维数。

\begin{theorem}[Poisson，二维波方程]\label{b2:thm:11-poisson-wave}
对二维光滑初值 \(f,g\)，唯一经典解为 \(u=\partial_tS_2(t)f+S_2(t)g\)，其中 \(t>0\) 时
\begin{equation}\label{b2:eq:11-two-dimensional}
 S_2(t)h(x)\coloneqq\dfrac1{2\pi c}
 \int_{|y-x|<ct}\dfrac{h(y)}{\sqrt{c^2t^2-|y-x|^2}}\,\mathrm{d}y.
\end{equation}
在 \(t=0\) 通过连续极限和时间导数恢复初值。
\end{theorem}
\begin{proof}
将上面的投影公式代入三维表示，先得到单位圆盘上 \(t/(2\pi)\) 乘加权积分。换元 \(y=x+ctz\) 后，二维体积因子为 \(c^2t^2\)，平方根带来 \(ct\)，得到\eqref{b2:eq:11-two-dimensional}。存在性和初值来自三维解，唯一性来自\cref{b2:prop:11-local-energy}。在单位圆盘的表达式里，权重固定且可积，光滑数据的时间求导有紧域上的共同控制；不应对移动圆盘上的奇异积分不加说明地直接使用边界求导公式。
\end{proof}

这称为二维波方程的\term[poisson-bodong-gongshi]{Poisson 公式}[Poisson's formula for the wave equation]，与调和函数的 Poisson 积分公式是不同的表示。若 \(f=0\)、\(g\geq0\) 光滑紧支撑且非零，对固定 \(x\) 和充分大的 \(t\)，整个初速度支集都位于积分圆盘内，而核严格为正，故 \(u(t,x)>0\)。二维波在经过以后留下逐渐减小的尾部。

\ProseParagraphBreak

任意奇维同样可由球平均给出有限阶微分表示。以下使用光滑数据以避免与维数有关的有限阶可微性记账；低正则数据随后在能量空间中单独处理。令 \(D_t=t^{-1}\partial_t\)，并约定空次幂为恒等算子，\((2m-1)!!=1\cdot3\cdots(2m-1)\)。

\begin{theorem}[奇维球面表示]\label{b2:thm:11-odd-wave}
设 \(n=2m+1\geq3\)，初值光滑。在 \(t>0\) 定义
\begin{equation}\label{b2:eq:11-odd-wave}
 S_n(t)h\coloneqq\dfrac1{(2m-1)!!}
 D_t^{m-1}\bigl(t^{2m-1}M_{ct}h\bigr).
\end{equation}
则 \(u=\partial_tS_n(t)f+S_n(t)g\) 为唯一经典解。初速度传播只依赖 \(\partial B_{ct}(x)\) 上有限阶的初值导数。
\end{theorem}
\begin{proof}
对 \(a(t,x)=M_{ct}h(x)\)，球平均恒等式成为
\[
 a_{tt}+\dfrac{2m}{t}a_t=c^2\Delta_xa.
\]
记 \(T_ma=D_t^{m-1}(t^{2m-1}a)\)。所需的微分恒等式为
\[
 (\partial_t^2-c^2\Delta_x)T_ma
 =T_m\left(a_{tt}+\dfrac{2m}{t}a_t-c^2\Delta_xa\right).
\]
可用乘积法则逐阶建立：\(D_t(t^ka)=kt^{k-2}a+t^{k-1}a_t\)，因而
\(T_m=t\prod_{j=1}^{m-1}(t\partial_t+2j+1)\)。把 \(t\partial_t\) 移过 \(t^k\) 时，按 \((t\partial_t)t^k=t^k(t\partial_t+k)\) 增加 \(k\)，两边即得到相同因子。比如 \(m=2\) 时 \(T_2a=3ta+t^2a_t\)，两边的时间项均为 \(t^2a_{ttt}+7ta_{tt}+8a_t\)。因此 \(S_n(t)h\) 满足波方程。

球平均为光滑偶函数，所以 \(a(t,x)=h(x)+O(t^2)\)，而此展开可在任意所需有限阶微分后继续使用。对常数项施加 \(T_m\) 得到 \((2m-1)!!\,t\cdot h\)，余项为 \(O(t^3)\)。于是 \(S_n(0)h=0\)、\(\partial_tS_n(0)h=h\)、\(\partial_t^2S_n(0)h=0\)。据此检验 \(u\) 的两项初值。唯一性仍由局部能量给出。最后，时间求导只改变球面半径及被积函数的有限阶导数，不把取值位置移入球内，故支集陈述成立。
\end{proof}

偶维 \(n=2m\geq2\) 从 \(n+1\) 维降下来。单位球面投影的同一计算给出
\begin{equation}\label{b2:eq:11-even-wave}
 S_n(t)h(x)=\dfrac2{(2m-1)!!\,\sigma_{2m}}
 D_t^{m-1}\left(t^{2m-1}
 \int_{|z|<1}\dfrac{h(x+ctz)}{\sqrt{1-|z|^2}}\,\mathrm{d}z\right).
\end{equation}
分母中的 \(\sigma_{2m}\) 是高一维单位球面的面积。公式的存在性、初值与唯一性由奇维结果和降维法得到；\(m=1\) 时恰好回到二维 Poisson 公式。\cref{b2:fig:11-kernel-support}将三维球面支集与二维圆盘支集并列，后者不能由前者简单删去一个坐标符号得到。
\begin{figure}[htbp]
\centering
\begin{tikzpicture}[x=1cm,y=1cm]
\begin{scope}[shift={(-2.3,0)}]
\draw[AnalysisBlue,very thick] (0,0) circle(1.25);
\fill (0,0) circle(1.2pt);
\draw[->] (0,0)--(1.2,0) node[midway,above]{\(ct\)};
\node[align=center] at (0,-1.9){三维：球面支集\\[2pt]图示为平面截面};
\end{scope}
\begin{scope}[shift={(2.3,0)}]
\fill[AnalysisBlue!16] (0,0) circle(1.25);
\draw[AnalysisBlue] (0,0) circle(1.25);
\fill (0,0) circle(1.2pt);
\draw[->] (0,0)--(1.2,0) node[midway,above]{\(ct\)};
\node[align=center] at (0,-1.9){二维：闭圆盘支集\\[2pt]内部具有非零核};
\end{scope}
\end{tikzpicture}

\smallskip
\begin{tikzpicture}[x=1cm,y=1cm]
\begin{scope}[shift={(-2.3,0)}]
\draw[AnalysisBlue,thick] (0,0) circle(1);
\draw[->] (-1.3,0)--(1.5,0) node[right]{\(z\)};
\draw[->] (0,-1.2)--(0,1.3) node[above]{\(\theta_3\)};
\draw[->,dashed] (.55,{sqrt(1-.55*.55)})--(.55,.04);
\draw[->,dashed] (.55,{-sqrt(1-.55*.55)})--(.55,-.04);
\fill (.55,{sqrt(1-.55*.55)}) circle(1.2pt);
\fill (.55,{-sqrt(1-.55*.55)}) circle(1.2pt);
\node at (0,-1.4){同一圆盘位置};
\end{scope}
\node[align=left] at (2.5,.25){上下半球各贡献一次\\
\(\mathrm d\sigma=(1-|z|^2)^{-1/2}\mathrm dz\)\\
靠近边缘，面积权重增大};
\end{tikzpicture}
\caption[二维与三维初速度核的支集]{二维与三维初速度核的支集示意，不按比例。上排左图只画球面的平面截线；右图为真实二维空间。下排为单位球面的过心截面，\((z,\theta_3)\) 表示球面方向坐标，\(z\) 为二维投影坐标，图中只画其中一个方向。垂直投影把上下两点送到同一圆盘位置；投影圆盘仅显示为一条直径。球面靠赤道趋于竖直，投影面积变小，故球面面积相对于投影面积的权重增大。比较的是固定正时间的核，不是任意初值的实际非零集合。}
\label{b2:fig:11-kernel-support}
\end{figure}


\subsection{基本解}
\label{b2:subsec:110302}

为避免混同两个对象，先固定正时间 \(t\)。把 \(S_n(t)\) 的卷积核记为 \(K_{n,c}(t,\cdot)\)，即对光滑测试函数 \(h\)，有 \(S_n(t)h=K_{n,c}(t,\cdot)*h\)。球面或圆盘公式及其有限阶导数，使该核成为一个紧支集分布。前两种显式形式为
\[
 K_{3,c}(t,\cdot)=\dfrac1{4\pi c^2t}\,\sigma_{ct},\quad
 K_{2,c}(t,x)=\dfrac{\mathbf1_{\{|x|<ct\}}}{2\pi c\sqrt{c^2t^2-|x|^2}},
\]
其中 \(\sigma_{ct}\) 为球面 \(\partial B_{ct}(0)\) 上的面积测度。两者对常数一的作用都为 \(t\)，它们是初速度核，不是单位质量的位移核。

\ProseParagraphBreak

再把这个核对负时间补零，记所得时空分布为 \(E_{n,c}\)。具体地，对时空测试函数 \(\varphi\)，规定
\[
 \langle E_{n,c},\varphi\rangle
 =\int_0^\infty\langle K_{n,c}(t,\cdot),\varphi(t,\cdot)\rangle\,\mathrm{d}t.
\]
零时刻 \(K(0)=0\)、\(\partial_tK(0)=\delta_0\) 的意义由对每个空间测试函数的极限给出。这里的分布配对仍采用双线性约定。

\begin{proposition}[延迟基本解]\label{b2:prop:11-retarded-fundamental}
时空分布 \(E_{n,c}\) 满足
\[
 (\partial_t^2-c^2\Delta_x)E_{n,c}=\delta_{(0,0)},\quad
 \operatorname{supp}E_{n,c}\subseteq\{\,(t,x):t\geq0,\ |x|\leq ct\,\}.
\]
对光滑外力和零初值，解为 \(\int_0^tS_n(t-s)F(s)\,\mathrm{d}s\)。
\end{proposition}
\begin{proof}
在 \(t>0\) 时，各维表示都满足齐次波方程。对时间积分分部两次，负半轴上的零延拓不产生位移跳跃，因为 \(K(0)=0\)；速度跳跃则为 \(\delta_0\)，给出时空原点的点质量。等价地，截断积分到 \(t\geq\varepsilon\)，其边界项为 \(\langle K_t(\varepsilon),\varphi(\varepsilon)\rangle-\langle K(\varepsilon),\varphi_t(\varepsilon)\rangle\)，在 \(\varepsilon\downarrow0\) 时趋于 \(\varphi(0,0)\)。传播支集来自固定时间核的球或球面支集；时间微分不会扩大闭支集。外力表达式的验证与\cref{b2:thm:09-duhamel}相同：第二次时间求导在积分上限产生 \(F(t)\)，其余项由齐次方程抵消。
\end{proof}

这一基本解因支集限制在未来而称为\term[yanchi-jibenjie]{延迟基本解}[retarded fundamental solution]。它的时空支集与某个给定初值所产生的解的支集不同，不能把基本解的整个传播区域误当成任意解必定非零的区域。

\ProseParagraphBreak

奇维 \(n\geq3\) 的固定时间核支撑在球面上；偶维核则具有球内尾部。后一个事实也可从\eqref{b2:eq:11-even-wave}精确读取：只在开球 \(|x|<ct\) 内考察其分布限制，换元后再对 \(t\) 求导得到光滑函数
\[
 K_{n,c}(t,x)=\dfrac{2(-1)^{(n-2)/2}}{(n-1)\sigma_n c}
 \bigl(c^2t^2-|x|^2\bigr)^{-(n-1)/2},\quad n\geq2\text{ 为偶数}.
\]
它在球内处处非零。对于 \(n\geq4\)，这个表达式在球面附近不再局部可积，因而只描述核在开球内部的限制；完整核必须按圆盘公式的分布微分解释，不能把这一个发散函数直接当作整个核积分。奇维无尾也不包括一维：\chapref{b2:ch:09}的初速度核是区间示性函数除以 \(2c\)。

\subsection{Fourier 方法}
\label{b2:subsec:110303}

球面表示显示几何，频率表示则统一各维数。沿用第一编的 \(2\pi\) 归一化，记 \(\rho(\xi)=2\pi c|\xi|\)。对 Schwartz 初值作空间变换，方程成为
\[
 \partial_t^2\widehat u(t,\xi)+\rho(\xi)^2\widehat u(t,\xi)=0,
\]
所以
\begin{equation}\label{b2:eq:11-fourier-wave}
 \widehat u(t,\xi)=\cos(t\rho)\cdot\widehat f(\xi)
             +\dfrac{\sin(t\rho)}{\rho}\cdot\widehat g(\xi).
\end{equation}
商在 \(\rho=0\) 处取值 \(t\)。它是可去奇点，不能因此对初速度附加零均值条件。

\ProseParagraphBreak

这些乘子也与前面的各维公式一致。对一个平面波 \(h(x)=e^{2\pi\mathrm{i}x\cdot\xi}\)，球面或圆盘算子仍给同一平面波乘一个标量；这个标量满足初值为 \((0,1)\) 的上述常微分方程，因此等于 \(\sin(t\rho)/\rho\)。Schwartz 反演把一般 \(h\) 写成平面波积分，快速下降控制表示中有限阶导数引入的频率幂，故可以交换积分和这些导数。于是两条构造路线在共同的光滑数据类上相同。

\begin{theorem}[任意维能量解]\label{b2:thm:11-fourier-energy}
对 \(f\in H^1(\mathbb R^n)\)、\(g\in L^2(\mathbb R^n)\)，\eqref{b2:eq:11-fourier-wave}通过 \(L^2\) 反演定义唯一的
\[
 u\in C(\mathbb R;H^1)\cap C^1(\mathbb R;L^2)
\]
分布解，满足相应初值。能量
\[
 E(t)=\dfrac12\bigl(\|u_t(t)\|_2^2+c^2\|\nabla u(t)\|_2^2\bigr)
\]
守恒，且 \(\|u(t)\|_2\leq\|f\|_2+|t|\|g\|_2\)。
\end{theorem}
\begin{proof}
由 \(|\sin(t\rho)/\rho|\leq|t|\)，公式定义 \(L^2\) 函数，并给出零阶估计。速度变换为 \(-\rho\sin(t\rho)\widehat f+\cos(t\rho)\widehat g\)。将其与 \(\rho\widehat u\) 合在一起，得到
\[
 \begin{pmatrix}\rho\widehat u(t)\\\widehat u_t(t)\end{pmatrix}
 =\begin{pmatrix}\cos(t\rho)&\sin(t\rho)\\-\sin(t\rho)&\cos(t\rho)\end{pmatrix}
 \begin{pmatrix}\rho\widehat f\\\widehat g\end{pmatrix}.
\]
右侧矩阵在每个频率保持复二维向量的欧氏平方范数。积分后使用\cref{b2:thm:04-unitary}与\cref{b2:thm:04-weak-derivative}，得到梯度平方可积及能量守恒。注意一阶梯度权重确为 \(4\pi^2|\xi|^2\)，没有高阶多重指标计数的问题。

时间连续性由频率逐点连续性和控制收敛得到：梯度、速度用固定的 \(\rho|\widehat f|+|\widehat g|\) 控制，位移在每个有界时间区间用 \(|\widehat f|+T|\widehat g|\) 控制。位移的差商也以速度的上述界控制，所以它在 \(L^2\) 中可微。先在有界频率截断上验证分布方程，再让截断半径增大；测试函数的导数属于 \(L^2\)，故范数收敛允许取极限。初值直接由时间连续性与公式恢复。

若另有同类解，取差并作空间磨光。\cref{b2:thm:09-energy-solution}证明末尾的时间微积分与全空间分部积分论证逐坐标相加，得到零数据磨光解的能量为零；再由初始位移为零得到解为零。磨光极限给出唯一性。
\end{proof}

外力 \(F\in L^1(0,T;L^2)\) 可用 Duhamel 积分加入。速度和梯度的频率旋转估计给出
\[
 \sqrt{2E(t)}\leq\sqrt{2E(0)}+\int_0^t\|F(s)\|_2\,\mathrm{d}s,
\]
而位移用 \(\|S_n(t-s)F(s)\|_2\leq(t-s)\|F(s)\|_2\) 控制。先对简单外力定义积分再取范数极限，重复\cref{b2:thm:09-energy-solution}的功积分极限论证，便得到非齐次能量恒等式。该构造不要求逐个空间点定义外力的值。

\subsection{支集性质}
\label{b2:subsec:110304}

Fourier 公式在所有频率上积分，有限传播在其中并不直观。支集信息可由共同光滑数据类上的几何表示传入能量极限：对支撑在闭集 \(A\) 内的数据先截断，再用半径 \(\varepsilon\) 的空间磨光核逼近。近似数据只扩大到 \(A\) 的 \(\varepsilon\) 邻域；经典有限传播与能量范数收敛使极限在 \(\operatorname{dist}(x,A)>ct\) 的开集上为零。因此，对每个固定时间，能量解仍满足同样的本性支集包含关系。

\ProseParagraphBreak

若 \(n\geq3\) 为奇数，光滑近似还满足壳层支集；初值支撑在 \(\overline B_R\) 时，先用半径 \(R+\varepsilon\) 的壳层，再令 \(\varepsilon\downarrow0\)，便得到\eqref{b2:eq:11-shell}的任意奇维能量版本。一般 \(H^1\) 数据在高维球面上未必具有公式所需的所有经典导数；壳层结论来自范数极限，不来自把某个任意代表直接限制到球面。

\ProseParagraphBreak

有限传播与频率局部化也须区别。若先把数据投影到一段频带，投影算子的空间卷积核通常有非紧支集，投影后的数据已经在远处留下尾部。下一节估计这种频带数据的逐点衰减，不能把投影造成的空间尾部归咎于波方程失去有限传播速度。几何支集陈述针对原始的完整演化，频带估计则针对先经一个非局部算子处理的数据。


% !TeX root = ../../Book2.tex
\OptionalSection{分散性观点}{b2:sec:1104}

能量守恒不排斥某些逐点振幅随时间减小：相同的平方能量可以分布在越来越大的区域上。为了看清波动传播中的振荡抵消，本节只证明三维、固定频带上的一个衰减估计。频带和空间正则性是结论的一部分；不能删掉频率限制而仍保持同一个 \(L^1\) 数据范数。

\subsection{\texorpdfstring{\(L^1\to L^\infty\)}{L¹→L∞} 衰减}
\label{b2:subsec:110401}

在 \(\mathbb R^3\) 上定义一支波传播算子
\[
 U_c(t)h\coloneqq\mathcal F_2^{-1}
 \left(e^{2\pi\mathrm{i}ct|\xi|}\widehat h(\xi)\right).
\]
它在 \(L^2\) 上等距；另一支以 \(-t\) 代替 \(t\)。取固定 \(\chi\in C_c^\infty((1/2,2))\)，用径向乘子
\[
 \widehat{P_Nh}(\xi)=\chi(|\xi|/N)\widehat h(\xi),\quad N>0,
\]
把频率局部化到大小为 \(N\) 的环带。这里称为\term[pinlv-jubuhua]{频率局部化}[frequency localization]，其空间作用是卷积，并不是把函数截断到某个空间区域。

\begin{theorem}[三维频带衰减]\label{b2:thm:11-localized-decay}
存在仅依赖 \(c\)、\(\chi\) 的常数 \(C\)，使任意 \(N>0\)、\(t\in\mathbb R\)、\(h\in L^1(\mathbb R^3)\) 满足
\begin{equation}\label{b2:eq:11-localized-decay}
 \|U_c(t)P_Nh\|_\infty
 \leq C N^3(1+N|t|)^{-1}\|h\|_1.
\end{equation}
此处由光滑紧频带的卷积核定义算子，因而不要求 \(h\in L^2\)。
\end{theorem}
\begin{proof}
先取 \(N=1\)。核为
\[
 K(t,x)=\int_{\mathbb R^3}e^{2\pi\mathrm{i}(x\cdot\xi+ct|\xi|)}
                       \chi(|\xi|)\,\mathrm{d}\xi.
\]
绝对值直接积分给出统一常数界，足以处理 \(|t|\leq1\)。令 \(r=|x|\)。若 \(r>0\)，对角变量积分，利用
\(\int_{\mathbb S^2}e^{2\pi\mathrm{i}r\rho\theta_1}\,\mathrm{d}\sigma
=2\sin(2\pi r\rho)/(r\rho)\)，得到
\[
 K(t,x)=\dfrac1{\mathrm{i}r}\bigl(I(ct+r)-I(ct-r)\bigr),\quad
 I(s)=\int_0^\infty\rho\chi(\rho)e^{2\pi\mathrm{i}s\rho}\,\mathrm{d}\rho.
\]
当 \(r\geq c|t|/2\) 时，\(I\) 的绝对值有统一界，因此 \(|K|\leq C/|t|\)。这个区域包括最难抵消的传播球面 \(r=c|t|\)。

若 \(r<c|t|/2\)，在角积分之前对径向变量操作。相位的径向导数为 \(ct+x\cdot\theta\)，绝对值至少为 \(c|t|/2\)。因 \(\rho^2\chi(\rho)\) 光滑且紧支撑在正半轴内部，两次分部积分没有边界项，给出
\[
 \left|\int_0^\infty\rho^2\chi(\rho)
 e^{2\pi\mathrm{i}\rho(ct+x\cdot\theta)}\,\mathrm{d}\rho\right|
 \leq C|t|^{-2}
\]
并且常数与方向无关。球面积分后仍有同样估计，也覆盖 \(r=0\)。两个区域合起来得到 \(\|K(t)\|_\infty\leq C(1+|t|)^{-1}\)。

一般 \(N\) 时换元 \(\xi=N\eta\)，有 \(K_N(t,x)=N^3K(Nt,Nx)\)。核的上确界遂满足\eqref{b2:eq:11-localized-decay}中的系数。最后对 \(K_N*h\) 使用绝对积分界，得到算子估计。
\end{proof}

当 \(N|t|\geq1\) 时，右端不超过 \(C|t|^{-1}N^2\|h\|_1\)。两个频率幂说明，全频率估计必须支付相应的空间正则性代价。例如对适当频带分解，若 \(h=\sum_jh_j\)、\(h_j\) 的频率支撑在大小为 \(2^j\) 的环带，且
\[
 \sum_j2^{2j}\|h_j\|_1<\infty,
\]
则用一个在每个环带上等于一的稍宽光滑截断，逐项应用核估计，得到
\[
 \left\|\sum_jU_c(t)h_j\right\|_\infty
 \leq\dfrac C{|t|}\sum_j2^{2j}\|h_j\|_1,\quad t\ne0.
\]
右侧使传播后的级数绝对一致收敛；若初始分解同时在分布意义下恢复 \(h\)，这个和就是完整演化。这里只给出一个明确的充分条件，不把加权频带和擅自替换为没有导数的 \(L^1\) 范数。

\subsection{色散}
\label{b2:subsec:110402}

在振荡积分中，不同方向的相位可以互相抵消；传播球面附近保留最大的贡献，而远离它的区域允许更多次分部积分。这种由相位和传播方向使振幅分散的现象，在本节称为\term[sesan]{色散}[dispersion]。波方程的关系式 \(\rho=2\pi c|\xi|\) 对径向频率是线性的，沿同一方向的不同频率以相同速度行进；因此它的衰减机制包含角向展开，不能直接照搬二次相位的热核或 Schrödinger 核计算。

\ProseParagraphBreak

频率代价不是上述证明的偶然损失。取非负非零 \(\chi\)，令 \(h=\mathcal F^{-1}(\chi(|\xi|))\)，再令 \(h_N(x)=N^3h(Nx)\)。这些都是 Schwartz 函数，且 \(\|h_N\|_1=\|h\|_1\)。在固定 \(t>0\) 和 \(|x|=ct\) 处，刚才的角积分公式给出
\[
 U_c(t)h_N(x)=\dfrac{N^2}{\mathrm{i}ct}
                    \bigl(I(2Nct)-I(0)\bigr).
\]
分部积分说明 \(I(2Nct)\to0\)，而 \(I(0)>0\)，故其模以 \(N^2\) 的量级增长。于是不存在对全部 Schwartz 数据成立的固定时间无导数 \(L^1\to L^\infty\) 界；更不能宣称完整波群具有 \(C|t|^{-1}\|h\|_1\) 的三维界。

\ProseParagraphBreak

另一方面，\(U_c(t)\) 在 \(L^2\) 中没有范数衰减。这与频带上的逐点衰减并不冲突：上确界下降时，波可以分布在更大的球壳上。三维壳层的面积按 \(t^2\) 增长，幅度按 \(t^{-1}\) 减小，恰与平方能量保持同一尺度相容。这只是对估计的几何解释；严格结论仍是上面带频带条件的核界与 Plancherel 等式。

\subsection{Strichartz 估计的动机}
\label{b2:subsec:110403}

非线性波方程中的外力依赖解本身，仅有每个时刻的能量界往往不够控制 Duhamel 积分。需要同时记录时间和空间的可积性，例如
\[
 \|u\|_{L^q_tL^r_x(I\times\mathbb R^n)}
 =\left(\int_I\left(\int_{\mathbb R^n}|u(t,x)|^r\,\mathrm{d}x\right)^{q/r}
 \,\mathrm{d}t\right)^{1/q},
\]
有限指数时按此定义，指数为无穷时改用本性上确界。能量给出某些 \(q=\infty\) 的控制，色散则提供随时间变化的空间映射界；连接两者的时空估计属于\term[strichartz-guji]{Strichartz 估计}[Strichartz estimates]的研究范围。

\ProseParagraphBreak

本节只证明了单个频带上的核估计。时间积分的端点、频带求和和数据的正则性都需要额外论证。陶哲轩的《Nonlinear Dispersive Equations: Local and Global Analysis》\cite{Tao2006Nonlinear}将此类时空估计用于非线性色散方程的局部与整体分析；本卷\appref{b2:app:E}将进一步讨论相应的时空估计。前面各节的球面表示、有限传播与能量解不依赖本选读节。



\begin{chaptersummary}
球面平均保留所有传播方向的对称信息，其径向方程中的维数系数决定了表示公式。三维乘以半径后化成一维波方程；奇维由有限阶半径微分得到球面表示，偶维经降维产生球内核。所有维数都有有限传播，强 Huygens 原理则只在奇维且维数至少为三时成立。

\ProseParagraphBreak

几何表示与 Fourier 公式在光滑数据上相合。前者显示支集，后者将每个频率变成一个保持能量的二维旋转，并自然处理 \(H^1\times L^2\) 初值。频率局部化后的三维核还具有 \(t^{-1}\) 衰减，但相应估计带有必要的频率代价；能量守恒、逐点衰减和时空可积性是不同的控制方式。后续热方程将出现另一种传播行为，不能只在这里的公式中更换时间因子就推断所有性质。
\end{chaptersummary}

\BookExercises

本组习题先考察球平均、径向解和正性，再以尾部渐近、有限时外力与尺度变化检验主线。最后两题使用选读部分的频带观点。

% 自拟：平均只记录二次型的迹，给出不能反推完整Hesse矩阵的例子。
\begin{exercise}\label{b2:ex:11-quadratic-mean}
设 \(A\) 为实对称 \(n\times n\) 矩阵，\(h(x)=x^{\mathsf T}Ax+b\cdot x+d\)。计算 \(M_rh(x)\)，并构造两个不同的二次多项式，使它们在原点的全部球面平均相同。
\end{exercise}
\begin{solution}
展开 \(h(x+r\theta)\)，一阶角项积分为零，二阶角项的平均为 \(r^2\operatorname{tr}A/n\)。因此
\[
 M_rh(x)=h(x)+\dfrac{r^2}{n}\operatorname{tr}A.
\]
当 \(n\geq2\) 时，\(h_1=0\) 与 \(h_2=x_1^2-x_2^2\) 的原点球面平均都为零，但二次型不同。故固定一个中心的球平均可以读取 Laplace 算子，却不能恢复全部 Hesse 矩阵。若 \(n=1\)，可取 \(h_1=0\)、\(h_2=x\)；关于原点的两点平均也都为零。
\end{solution}

% 自拟：三维径向解的原点可去奇性与一维奇延拓相接。
\begin{exercise}\label{b2:ex:11-radial-reduction}
设三维光滑径向初值为 \(f(x)=F(|x|)\)、\(g(x)=G(|x|)\)。令 \(v(t,r)=r\cdot u(t,r)\)，其中径向解以半径为变量。用奇延拓求 \(r>0\) 的解，并计算原点的极限值。
\end{exercise}
\begin{solution}
径向 Laplace 算子为 \(u_{rr}+2u_r/r\)，所以 \(v_{tt}=c^2v_{rr}\)。初值在正轴为 \(rF(r)\)、\(rG(r)\)，在负轴取奇延拓 \(H(r)=rF(|r|)\)、\(J(r)=rG(|r|)\)。原始径向数据的光滑性保证延拓在原点光滑。于是
\[
 u(t,r)=\dfrac{H(r+ct)+H(r-ct)}{2r}
       +\dfrac1{2cr}\int_{r-ct}^{r+ct}J(s)\,\mathrm{d}s.
\]
分子在 \(r=0\) 为零。对 \(r\) 求导后令 \(r=0\)，得到
\[
 u(t,0)=F(ct)+ctF'(ct)+tG(ct),\quad t>0.
\]
这也等于球平均公式在中心的值。分母中的 \(r\) 因奇性而产生可去奇点，不能把它理解为原点真实出现的无穷大。
\end{solution}

% 自拟：三维正位移不保持正性，与一维平均公式形成边界比较。
\begin{exercise}\label{b2:ex:11-positivity-failure}
证明三维齐次波方程不保持初始位移的非负性，即使初速度为零也如此。可取 \(f(x)=e^{-|x|^2}\)。解释这与一维 \(g=0\) 时的 d'Alembert 公式有何不同。
\end{exercise}
\begin{solution}
在原点使用径向公式，得到
\[
 u(t,0)=\bigl(1-2c^2t^2\bigr)e^{-c^2t^2}.
\]
当 \(ct>1/\sqrt2\) 时它为负，而初始位移处处为正。三维位移项包含球面法向导数，这一项不保持符号。一维零初速度解仅是初始位移两个平移的算术平均，因此在前后任意时间都保持位移非负；不能从三维初速度核为正误推出其时间导数也保持正性。
\end{solution}

% 自拟：二维尾部的两项渐近，增加总质量和二阶矩的辨析。
\begin{exercise}\label{b2:ex:11-two-dimensional-tail}
设二维 \(f=0\)、\(g\in C_c^\infty(\mathbb R^2)\)。证明对固定 \(x\)，当 \(t\to\infty\) 时
\[
 u(t,x)=\dfrac1{2\pi c^2t}\int g(y)\,\mathrm{d}y
 +\dfrac1{4\pi c^4t^3}\int|x-y|^2g(y)\,\mathrm{d}y+O(t^{-5}).
\]
并说明该展开在 \(x\) 的每个固定紧集上一致。
\end{exercise}
\begin{solution}
对充分大的 \(t\)，数据支集全部位于积分圆盘内。在固定的 \(x\) 紧集及 \(g\) 支集上，\(|x-y|\leq M\)。取 \(ct\geq2M\)，使用
\[
 (1-z)^{-1/2}=1+\dfrac z2+O(z^2),\quad 0\leq z\leq1/4,
\]
并令 \(z=|x-y|^2/(c^2t^2)\)。把它代入\eqref{b2:eq:11-two-dimensional}，得到前两项；余项被 \(Ct^{-5}|x-y|^4|g(y)|\) 控制，其积分有限且在该紧集上一致。若 \(\int g=0\)，首项消失，但通常仍有 \(t^{-3}\) 的尾部，零总初速度并不自动恢复 Huygens 原理。
\end{solution}

% 自拟：有限持续外力给出随源持续时间加厚的三维壳层。
\begin{exercise}\label{b2:ex:11-forcing-shell}
设三维波方程初值为零，光滑外力支撑在 \([0,T_0]\times\overline{B_R(0)}\) 内。证明当 \(t>T_0\) 时
\[
 \operatorname{supp}u(t)\subseteq
 \{\,x:\max\{0,c(t-T_0)-R\}\leq|x|\leq ct+R\,\}.
\]
为什么不能把内半径直接写成 \(ct-R\)？
\end{exercise}
\begin{solution}
时刻 \(s\) 的响应为 \(S_3(t-s)F(s)\)，支撑在由 \(|x|\leq c(t-s)+R\)、\(|x|\geq c(t-s)-R\) 描述的壳层内。对 \(0\leq s\leq T_0\) 取并，外半径至多为 \(ct+R\)，内半径至少为 \(c(t-T_0)-R\)，负值时改取零。对这些响应积分不会扩大到这个闭并集之外。临近 \(T_0\) 产生的波只传播了约 \(t-T_0\) 的时间，其后缘仍更靠近原点，因此不能使用只适合零时刻单次扰动的内半径。
\end{solution}

% 自拟：空间时间同比缩放下能量不变与零阶范数不同尺度。
\begin{exercise}\label{b2:ex:11-energy-scaling}
对 \(n\) 维齐次解定义
\[
 u_\lambda(t,x)=\lambda^{n/2-1}u(\lambda t,\lambda x),\quad\lambda>0.
\]
求其初值，证明它仍具有相同波速且总能量与 \(u\) 相同，并求 \(\|u_\lambda(0)\|_2\) 的变化。
\end{exercise}
\begin{solution}
时间二阶导数和空间二阶导数都多出 \(\lambda^2\)，所以波速不变。初值为 \(f_\lambda(x)=\lambda^{n/2-1}f(\lambda x)\)、\(g_\lambda(x)=\lambda^{n/2}g(\lambda x)\)。速度和梯度均有振幅因子 \(\lambda^{n/2}\)，其平方积分的变量代换再贡献 \(\lambda^{-n}\)，故能量与原解相同。零阶范数却满足 \(\|f_\lambda\|_2=\lambda^{-1}\|f\|_2\)。因此保持能量的缩放不保持非齐次 \(H^1\) 范数，低频的零阶控制不能由能量缩放独自决定。
\end{solution}

% 自拟：初速度传播的频带损失比半波少一阶，计算传播球面下界。
\begin{optionalexercise}\label{b2:ex:11-velocity-band-loss}
在三维取非负非零 \(\chi\in C_c^\infty((1/2,2))\)。证明
\[
 \|S_3(t)P_Nh\|_\infty\leq C N^2(1+N|t|)^{-1}\|h\|_1.
\]
再在固定 \(t>0\)、\(|x|=ct\) 上计算这个乘子的卷积核的首项，证明其量级为 \(N/t\)。
\end{optionalexercise}
\begin{solution}
乘子为 \(\sin(2\pi ct|\xi|)\chi(|\xi|/N)/(2\pi c|\xi|)\)。在支集上分母可写成 \(N\) 乘一个光滑有界因子，将正弦拆成两支指数后，\cref{b2:thm:11-localized-decay}给出所求界。对核做角积分，在 \(r=|x|=ct\) 时得到
\[
 \dfrac1{2\pi c^2t}\int_0^\infty
 \chi(\rho/N)\bigl(1-\cos(4\pi ct\rho)\bigr)\,\mathrm{d}\rho.
\]
第一项为 \(N(2\pi c^2t)^{-1}\int\chi\)，严格为正；第二项换元后对光滑紧支集函数分部积分，成为 \(o(N)\)。因此核的首项具有 \(N/t\) 量级。初速度传播比单支波传播少一个频率幂，来自 \(1/|\xi|\) 因子，但仍不能获得固定时间的全频无正则性 \(L^1\to L^\infty\) 界。
\end{solution}

% 自拟：同时有限空间支集和有限频率支集只能为零，以整指数级数证明。
\begin{optionalexercise}\label{b2:ex:11-bandlimited-support}
设 \(h\in L^2(\mathbb R^n)\) 及其 Fourier 变换都具有紧支集。只利用 Fourier 反演和指数级数证明 \(h=0\)，不调用复分析中的恒等定理。由此解释频率局部化后为何不应继续假设数据具有原来的紧空间支集。
\end{optionalexercise}
\begin{solution}
频率支集有限且 \(\widehat h\in L^2\)，故由 Cauchy--Schwarz 有 \(\widehat h\in L^1\)。反演给出连续代表。固定其余坐标 \(x'\)，把第一坐标记为 \(s\)。若频率支集包含在 \(|\xi|\leq B\) 内，则对任意实数 \(s_0\)，
\[
 h(s,x')=\sum_{k=0}^\infty\dfrac{(s-s_0)^k}{k!}
 \int(2\pi\mathrm{i}\xi_1)^k\widehat h(\xi)
       e^{2\pi\mathrm{i}(s_0\xi_1+x'\cdot\xi')}\,\mathrm{d}\xi.
\]
各项绝对值的和不超过 \(e^{2\pi B|s-s_0|}\|\widehat h\|_1\)，所以逐项积分有效且级数对所有 \(s\) 收敛。空间紧支集使连续代表在充分大的 \(s_0\) 附近恒为零，因此所有这些导数系数都为零；级数便说明沿整条直线为零。让 \(x'\) 任意，得到 \(h=0\)。因此一个非零紧支集数据经非平凡有限频带截断后，若结果非零，就必然失去紧空间支集。
\end{solution}

